% v2026a; generated by R/47_OA_Examples.R (restricted stage); manifesto text from the Manifesto Project corpus
{\footnotesize
\begin{longtable}{p{0.96\linewidth}}
\caption{Worked examples: liberalism dimensions}\label{tab:oa9c} \\
\toprule
\textbf{Political liberalism} --- Liberal Party (United Kingdom), election May 1979; manifesto 51420\_197905, chunk 1. Scores: Sonnet 9, GPT-4.1-mini 9, Gemini 9. \\
\quad \emph{Quote (Sonnet 4.6):} ``grossly undemocratic voting system'' \\
\quad \emph{Context:} Britain has a grossly undemocratic voting system, over-centralised government and an ineffective Parliament. Piecemeal changes have failed to introduce the necessary constitutional reforms. \\
\quad \emph{Justification (liberalism block):} The manifesto is strongly pro-liberal across all four dimensions: it champions proportional representation, a Bill of Rights, freedom of information, and democratic reform (political); equal opportunities, minority rights, and anti-discrimination legislation (social); free markets, competition, reduced taxation, and opposition to further nationalisation (economic); and references to interest rate management and market-enabling frameworks (financial). \\
\midrule
\textbf{Social liberalism} --- Green Party of England and Wales (United Kingdom), election May 2015; manifesto 51110\_201505, chunk 1. Scores: Sonnet 9, GPT-4.1-mini 9, Gemini 9. \\
\quad \emph{Quote (Gemini 3.5 Flash):} ``Ensure respect for everyone whatever their ethnicity, gender, age, religious belief'' \\
\quad \emph{Context:} Ensure respect for everyone whatever their ethnicity, gender, age, religious belief or non-belief, sexual orientation, class, size, disability or other status. \\
\quad \emph{Justification (liberalism block):} The manifesto strongly supports political and social liberalism through democratic reforms and minority rights, but rejects economic and financial liberalism in favor of state intervention and public ownership. \\
\midrule
\textbf{Economic liberalism} --- Conservative Party (United Kingdom), election April 1992; manifesto 51620\_199204, chunk 1. Scores: Sonnet 9, GPT-4.1-mini 9, Gemini 9. \\
\quad \emph{Quote (Gemini 3.5 Flash):} ``Competition and private ownership are the most'' \\
\quad \emph{Context:} PRIVATISATION Competition and private ownership are the most powerful engines of economic efficiency, innovation and choice. \\
\quad \emph{Justification (liberalism block):} The text strongly advocates classical liberal principles, including free enterprise, privatization, deregulation, and the rule of law. \\
\midrule
\textbf{Financial-market liberalism} --- Conservative Party (United Kingdom), election June 1987; manifesto 51620\_198706, chunk 1. Scores: Sonnet 8, GPT-4.1-mini 8, Gemini 8. \\
\quad \emph{Quote (Gemini 3.5 Flash):} ``foremost financial centre in the world'' \\
\quad \emph{Context:} by enabling the City of London to become the foremost financial centre in the world, by keeping down \\
\quad \emph{Justification (liberalism block):} The manifesto strongly champions classical liberal principles, including free enterprise, privatization, lower taxation, and the expansion of financial and property ownership. \\
\midrule
\textbf{Liberalism (overall)} --- Labour Party (United Kingdom), election May 2015; manifesto 51320\_201505, chunk 2. Scores: Sonnet 8, GPT-4.1-mini 8, Gemini 8. \\
\quad \emph{Quote (Sonnet 4.6):} ``individual freedom and democracy at the heart'' \\
\quad \emph{Context:} We will support human rights, always putting individual freedom and democracy at the heart of our foreign policy. \\
\quad \emph{Justification (liberalism block):} The text is strongly pro-liberal across political (democracy, rule of law, human rights in foreign policy), social (LGBT rights envoy, women's rights, religious freedom), and economic (single market completion, capital inflows, free trade) dimensions, with moderate support for open financial markets through EU membership and capital flows. \\
\bottomrule
\multicolumn{1}{p{0.96\linewidth}}{\scriptsize For each dimension the chunk of an English-language manifesto (United Kingdom, Ireland, Australia, New Zealand, United States) on which all three models give a verified score of at least 7 with a quote that passes the quote check and agree most closely (smallest range; ties: higher mean, then chunk id). Evidence shown: the model with the shortest context. Quote, context and justification as returned by the model. Manifesto text: Manifesto Project corpus \citep{lehmann2026manifesto}, identified by the manifesto id.} \\
\end{longtable}
}
